\section{Conclusion} 
\label{sec_conclusion}

In this paper, we addressed the critical vulnerability of Federated Learning (FL) to model poisoning attacks, a challenge significantly amplified in realistic non-IID data environments. Existing defense mechanisms are often passive, computationally expensive, or struggle to distinguish between malicious updates and benign statistical diversity, leading to high false-positive rates. To overcome these limitations, we introduced P4P (Probe for Poisoning), a proactive, lightweight, and multi-layer defense framework.

The core innovation of P4P is the probe vector, generated from the global model's own historical trajectory, which serves as a common reference to actively query clients. By evaluating the directional consistency of each client's update via a simple cosine similarity response, P4P transforms a high-dimensional analysis problem into a lightweight, one-dimensional anomaly detection task. This proactive probing is integrated into a multi-layer defense pipeline, combining L2-norm filtering, ensemble anomaly detection, and a temporal suspicion scoring system to effectively counter a wide spectrum of attacks.

Our comprehensive experiments demonstrate that P4P significantly enhances the robustness of FL systems. It consistently maintains high model accuracy and stable convergence under various sophisticated poisoning attacks, including sign-flipping, label-flipping, and backdoors, where undefended models collapse. Crucially, P4P demonstrates strong resilience in highly heterogeneous non-IID settings, successfully identifying adversarial behavior without penalizing benign clients. Furthermore, its design ensures negligible computational and communication overhead, making it a practical and scalable solution for real-world, resource-constrained environments like IoT-based Intrusion Detection Systems. P4P represents a significant step towards building more secure and trustworthy decentralized machine learning systems.

\section*{Acknowledgement}

This research is funded by Vietnam National University HoChiMinh City (VNU-HCM) under grant number NCM2025-26-01.

% This work was supported by the VNUHCM–University of Information Technology's Scientific Research Support Fund.
